\chapter{Introducción}
\label{cap:introduccion}

En el Informe 2 se reconstruyó, mediante ingeniería inversa, el funcionamiento de BeeTracer,
la plataforma de Bilix Ingeniería que digitaliza el desconsolidado de contenedores en
terminales extraportuarios. Ese trabajo usó el Lenguaje Unificado de Modelado (UML) y el
enfoque orientado a objetos de \textcite{larman2003}: se identificaron los actores, los casos de
uso, las interacciones entre objetos y un diagrama de clases que explica la estructura del
sistema.

El presente informe corresponde a la tercera entrega de la asignatura Modelamiento de Software
y aborda el mismo sistema desde otro paradigma: el \textbf{análisis y diseño estructurado}. En
lugar de preguntar qué objetos colaboran, se pregunta qué datos entran al sistema, cómo se
transforman y dónde se almacenan. Para ello se sigue el análisis estructurado moderno de
\textcite{yourdon1989} y \textcite{demarco1979}, y el diseño estructurado de
\textcite{yourdon1979} y \textcite{pagejones1988}. Las fuentes de información son las mismas
del informe anterior: la charla técnica del equipo de BeeTracer
\parencite{charlabeetracer2026}, la información pública de la plataforma \parencite{bilix_sf}
y el propio modelo UML, que aquí se usa como punto de partida para asegurar la coherencia
entre ambas entregas.

El documento se organiza en dos partes, de acuerdo con la pauta de la asignatura:

\begin{itemize}
    \item \textbf{Parte 1, Análisis estructurado.} Comprende el \emph{modelo ambiental}, que
    delimita el sistema frente a su entorno mediante la declaración de propósitos, el
    diagrama de contexto y la lista de acontecimientos; y el \emph{modelo de comportamiento},
    que describe su funcionamiento interno mediante los diagramas de flujo de datos (niveles
    0, 1 y 2), el diagrama entidad-relación, las mini especificaciones de los procesos
    elementales y el diccionario de datos.
    \item \textbf{Parte 2, Diseño estructurado.} Traduce el modelo anterior a una
    arquitectura implementable: el diagrama de estructura, obtenido de los diagramas de flujo
    de datos mediante el análisis de transacción y de transformación; el modelo relacional,
    derivado del diagrama entidad-relación; y la especificación de cada módulo.
\end{itemize}

Siguiendo a \textcite{yourdon1989}, el análisis se plantea como un \emph{modelo esencial}: se
supone una tecnología perfecta, por lo que no se modelan aspectos como la falta de cobertura
en el patio, la compresión de imágenes o el canal por el que viaja un correo. Esos aspectos,
centrales en BeeTracer, no se pierden: reaparecen en la Parte 2 como módulos de
implementación (sincronización diferida, compresión de fotografías, envío de correo), que es
donde corresponde tratarlos. Al igual que en el informe anterior, los procesos, datos y
módulos presentados son una propuesta inferida, coherente con el comportamiento observado,
pero no necesariamente idéntica a la implementación real de BeeTracer.
